% ============================================================================
% AERIS Paper - Section 2: Related Work
% Target Journal: Applied Intelligence (APIN) - Springer
% ============================================================================

\section{Related Work}
\label{sec:related}

This section reviews existing WSN routing protocols with emphasis on the
latency-energy-reliability trade-off, positioning AERIS within the broader
research landscape.

\subsection{Classical Clustering Protocols}

Clustering-based routing has been the dominant paradigm for energy-efficient
WSN communication since the introduction of LEACH \cite{heinzelman2000leach}.
We categorize classical protocols by their transmission structure and analyze
their latency characteristics.

\subsubsection{Cluster-Based Protocols (LEACH, HEED)}

LEACH \cite{heinzelman2000leach} introduced the concept of rotating cluster
heads to distribute energy consumption. In each round, nodes self-elect as
cluster heads with probability $p$, and non-cluster-head nodes join the
nearest cluster head. Data transmission follows a two-phase model:
intra-cluster aggregation followed by direct cluster-head-to-base-station
transmission.

\textbf{Latency Analysis}: LEACH achieves $O(1)$ transmission latency since
all cluster heads transmit directly to the base station in parallel. However,
this direct transmission incurs high energy consumption for distant cluster
heads, and the random cluster head selection leads to suboptimal clustering
in heterogeneous deployments.

HEED \cite{younis2004heed} improved upon LEACH by incorporating residual
energy and communication cost into cluster head selection. The protocol
iteratively increases cluster head probability based on node energy levels,
achieving more balanced energy consumption. HEED maintains $O(1)$ latency
through its single-hop cluster-head-to-base-station model but introduces
additional setup overhead.

\textbf{Limitation}: Both LEACH and HEED exhibit PDR degradation at scale
due to increased transmission distances. Our experiments show LEACH PDR
drops to 98.7\% at 500 nodes (Section~\ref{sec:results}).

\subsubsection{Chain-Based Protocols (PEGASIS)}

PEGASIS \cite{lindsey2002pegasis} pioneered chain-based data aggregation,
where nodes form a linear chain and data is passed sequentially toward a
designated leader node. Each node receives data from one neighbor, fuses
it with its own data, and transmits to the next neighbor in the chain.

\textbf{Energy Efficiency}: PEGASIS achieves optimal energy efficiency by
minimizing transmission distances---each node only communicates with its
immediate neighbors. Our experiments confirm PEGASIS consumes 50\% less
energy than LEACH (41.9mJ vs 100.7mJ at 100 nodes).

\textbf{Latency Analysis}: The sequential transmission model introduces
$O(n)$ latency, where $n$ is the number of nodes. For a network with $n$
nodes, data must traverse an average of $n/2$ hops before reaching the
base station. At 500 nodes, this translates to approximately 2500ms
end-to-end delay---a critical limitation for real-time applications.

\textbf{Robustness}: Chain-based protocols are vulnerable to node failures.
A single node failure can partition the chain, requiring complete chain
reconstruction with $O(n^2)$ complexity.

\subsubsection{Summary of Classical Protocol Trade-offs}

Table~\ref{tab:classical_comparison} summarizes the trade-offs among
classical protocols.

\begin{table}[htbp]
\centering
\caption{Comparison of Classical WSN Routing Protocols}
\label{tab:classical_comparison}
\begin{tabular}{@{}lccccc@{}}
\toprule
Protocol & Latency & Energy & PDR & Robustness & Complexity \\
\midrule
LEACH & $O(1)$ & High & Degrades & Medium & $O(n)$ \\
HEED & $O(1)$ & Medium & Stable & Medium & $O(n \log n)$ \\
PEGASIS & $O(n)$ & \textbf{Low} & \textbf{High} & Low & $O(n^2)$ \\
\bottomrule
\end{tabular}
\end{table}

\subsection{Machine Learning-Based Routing}

Recent advances in machine learning have motivated numerous ML-based WSN
routing protocols that adapt to dynamic network conditions
\cite{ren2024mefi, okine2024madrl}.

\subsubsection{Reinforcement Learning Approaches}

Deep Q-Networks (DQN) have been applied to WSN routing to learn optimal
relay selection policies \cite{okine2024madrl}. Multi-Agent Deep
Reinforcement Learning (MADRL) frameworks enable distributed decision-making
where each node acts as an independent agent learning cooperative routing
strategies.

MeFi \cite{ren2024mefi} employs GRU-based mean field reinforcement learning
for cooperative routing, demonstrating improved adaptation to network
dynamics through centralized training with decentralized execution.

\textbf{Computational Barriers}: Despite promising results, ML-based
approaches face significant deployment challenges on resource-constrained
WSN nodes:

\begin{itemize}
    \item \textbf{Inference latency}: Neural network forward propagation
    requires 50--600ms on microcontroller-class processors, exceeding
    real-time requirements for industrial monitoring ($<$100ms)
    \cite{kumar2020tinyml}.

    \item \textbf{Memory footprint}: Even compact LSTM/GRU models require
    700KB--2MB for weight storage, far exceeding available RAM on commodity
    sensor nodes (TelosB: 10KB, CC2650: 20KB).

    \item \textbf{Training overhead}: RL algorithms typically require
    thousands of training episodes (8--96 hours on GPU infrastructure),
    making them impractical for dynamic deployments where conditions
    change post-installation.
\end{itemize}

Table~\ref{tab:ml_comparison} quantifies these computational requirements.

\begin{table}[htbp]
\centering
\caption{Computational Requirements of ML-Based Routing Protocols}
\label{tab:ml_comparison}
\begin{tabular}{@{}lcccc@{}}
\toprule
Method & Decision Time & Memory & Training & Hardware \\
\midrule
LSTM-routing & 50--80ms & 700KB & 16h & 512KB+ RAM \\
MeFi (GRU) & $\sim$600ms & 2MB & 48h & 1MB+ RAM \\
MADRL (DQN) & $\sim$500ms & 3.5MB & 96h & 2MB+ RAM \\
\textbf{AERIS} & $<$\textbf{10ms} & \textbf{23KB} & \textbf{0h} & \textbf{10KB+ RAM} \\
\bottomrule
\end{tabular}
\end{table}

\subsection{Environment-Aware Routing}

Environment-aware routing protocols incorporate physical environmental
factors (temperature, humidity, interference) into routing decisions
\cite{sun2021environment, wang2022adaptive}.

\subsubsection{Link Quality Prediction}

Several works have demonstrated strong correlations between environmental
conditions and wireless link quality. Temperature variations affect
transceiver noise figures, while humidity influences signal propagation
characteristics \cite{boano2010impact}. These findings motivate
environment-adaptive protocols that adjust transmission parameters based
on sensed conditions.

\subsubsection{Adaptive Transmission Strategies}

Cross-layer protocols integrate physical layer information (RSSI, LQI)
into routing decisions. Adaptive power control mechanisms adjust
transmission power based on estimated channel conditions, balancing
energy consumption against delivery reliability.

\textbf{Limitation}: Existing environment-aware approaches typically
focus on link-level adaptation without addressing the fundamental
latency-energy trade-off at the network level. Most retain either
$O(1)$ (cluster-based) or $O(n)$ (chain-based) latency characteristics
inherited from their underlying routing structures.

\subsection{Research Gap and AERIS Positioning}

Based on our analysis, we identify the following research gap:

\textbf{Gap Statement}: No existing lightweight protocol ($<$50KB memory,
$<$50ms decision time) simultaneously achieves:
\begin{enumerate}
    \item Sub-linear transmission latency: $O(\log n)$ or better
    \item High reliability at scale: $>$99\% PDR at 500 nodes
    \item Zero training requirement: Immediate deployment capability
    \item Commodity hardware compatibility: Deployable on 10KB RAM nodes
\end{enumerate}

Figure~\ref{fig:design_space} illustrates the protocol design space,
showing how AERIS fills the gap between energy-optimal (PEGASIS) and
latency-optimal (LEACH) approaches.

\begin{figure}[htbp]
\centering
% [Design space diagram showing LEACH, PEGASIS, HEED, ML methods, and AERIS positions]
\includegraphics[width=0.9\columnwidth]{fig_design_space}
\caption{WSN Protocol Design Space. AERIS occupies the previously unfilled
region combining low latency ($O(\log n)$), high PDR (100\%), and lightweight
implementation (23KB). Dashed lines indicate real-time latency threshold
(500ms) and commodity hardware memory limit (50KB).}
\label{fig:design_space}
\end{figure}

\textbf{AERIS Positioning}: AERIS addresses this gap through hierarchical
routing that achieves $O(\log n)$ latency while maintaining deterministic,
lightweight operation suitable for commodity WSN hardware. Unlike ML-based
approaches, AERIS requires no training and provides interpretable decision
logic. Unlike PEGASIS, AERIS sacrifices some energy efficiency (2$\times$
higher consumption) to achieve 96\% latency reduction.

Table~\ref{tab:positioning} summarizes AERIS positioning relative to
existing approaches.

\begin{table}[htbp]
\centering
\caption{AERIS Positioning in the Protocol Design Space}
\label{tab:positioning}
\begin{tabular}{@{}lccccc@{}}
\toprule
Approach & Latency & Energy & PDR & Training & Memory \\
\midrule
LEACH & \checkmark & $\times$ & $\times$ & \checkmark & \checkmark \\
PEGASIS & $\times$ & \checkmark & \checkmark & \checkmark & \checkmark \\
HEED & \checkmark & $\triangle$ & $\triangle$ & \checkmark & \checkmark \\
ML-based & $\times$ & $\triangle$ & \checkmark & $\times$ & $\times$ \\
\textbf{AERIS} & \checkmark & $\triangle$ & \checkmark & \checkmark & \checkmark \\
\bottomrule
\end{tabular}
\begin{tablenotes}
\small
\item \checkmark: Satisfies requirement; $\times$: Does not satisfy;
$\triangle$: Partially satisfies
\end{tablenotes}
\end{table}
